% =====================================================================
% Ch1 report — Dense networks on Fashion MNIST (1 A4)
% YOU (the student) write ALL prose. This file only provides structure.
%
% WRITE:
%   - Hypothesis: a real, testable question about how hyperparameters
%     INTERACT (e.g. width vs depth on simple 28x28 data) — not "tuning helps".
%   - Experiment design: the six sweeps (units 4x4, depth 2/3/4, learning rate,
%     optimizer, batch size, epochs), the ranges, and why powers of 2.
%   - Results: reference Fig.~\ref{fig:units} (key figure), optionally depth/LR.
%   - Conclusions.
% REFLECT ON:
%   - Did width help more than depth here, and WHY for a simple dataset?
%   - Epochs: upside (learning) vs downside (overfitting, compute).
%   - Powers of 2: upside (hardware alignment) vs downside (coarse search).
%   - Why the learning-rate sweet spot sits where it does; why 1e-5 barely learns.
%   - Did results match the hypothesis; what would you test next?
% =====================================================================
\documentclass[11pt,a4paper]{article}
\usepackage[margin=1.8cm]{geometry}
\usepackage{graphicx}
\usepackage{float}
\usepackage{caption}
\usepackage{booktabs}
\usepackage{fancyhdr}
\usepackage{parskip}
\setlength{\parskip}{2pt}
\usepackage{titlesec}
\titlespacing*{\section}{0pt}{4pt}{2pt}
\captionsetup{font=small,skip=3pt}
\setlength{\intextsep}{5pt}
\setlength{\textfloatsep}{5pt}

\newcommand{\studentname}{Vincent Brand}
\newcommand{\studentnumber}{1447806}
\newcommand{\course}{Machine Learning -- MADS, HAN}
\newcommand{\reportdate}{2026}
\newcommand{\repourl}{https://github.com/vincentbrand/MADS-ML-Vincent}

\pagestyle{fancy}
\fancyhf{}
\lhead{\studentname\ (\studentnumber)}
\rhead{\course}
\cfoot{\thepage}
\renewcommand{\headrulewidth}{0.4pt}
\sloppy

% include a figure if it exists, otherwise show a placeholder box.
% optional first arg = width (default \linewidth) so figures fit on 1 A4.
\newcommand{\figorbox}[2][\linewidth]{%
  \IfFileExists{#2}{\includegraphics[width=#1]{#2}}%
  {\fbox{\parbox[c][3.5cm][c]{0.98\linewidth}{\centering\ttfamily
   missing figure: #2\\[2pt] run the notebook + export step to generate it}}}}

\begin{document}

\begin{center}
  {\Large\bfseries Hyperparameter experiments on Fashion MNIST}\\[3pt]
  \studentname\ \textbar\ \studentnumber\ \textbar\ \course\\
  \reportdate\ \textbar\ \repourl
\end{center}

\section{Hypothesis}
Fashion-MNIST is a small ($28\times28$, single-channel) 10-class problem, so I
expect a dense network's performance to be limited mainly by the \emph{width} of
its hidden layers rather than its \emph{depth}. My hypothesis is that on this
simple data, widening the layers lowers test loss more than stacking extra
layers, and that the \emph{first} hidden layer --- the one that sees the raw
784-pixel input --- drives most of that gain. I also expect a clear
learning-rate sweet spot and diminishing returns past a few epochs.

\section{Experiment design}
Every run uses the provided three-\texttt{Linear}-layer network
(\texttt{CrossEntropyLoss}, Adam, \texttt{ReduceLROnPlateau}); I report the final
test loss. From a $128\times128$ / lr$=10^{-3}$ / batch$=64$ /
5-epoch baseline I changed one factor at a time across six sweeps: a $4\times4$
width grid over units1$\times$units2 $\in\{32,64,128,256\}$
(Fig.~\ref{fig:units}); depth (2/3/4 layers at 64/128/256 units); learning rate
($10^{-2}$--$10^{-5}$); optimizer (SGD/Adam/AdamW/RMSprop); batch size (8--128);
and epochs (3/5/10) (Fig.~\ref{fig:sweeps}). I used powers of two because they
align with how memory and compute tiles are allocated on the hardware; the
trade-off is a coarse grid that can step over a better value in between.

\section{Results}
As shown in Fig.~\ref{fig:units}, wider networks consistently reach lower test
loss: the best two-layer model is $256\times256$ ($0.334$), the worst
$32\times32$ ($0.405$). Averaged across the grid, widening the \emph{first} layer
from 32 to 256 units cuts loss by $\approx0.040$ but the \emph{second} layer by
only $\approx0.018$ --- the first hidden layer matters about twice as much, as
expected since it must compress the full 784-pixel input. Depth did not help: at
256 units the two-layer net ($0.331$) beat the three- and four-layer ones
($0.339$, $0.345$), and at 128/64 units it was flat. The sweeps
(Fig.~\ref{fig:sweeps}) give a U-shaped learning-rate curve minimised near
$10^{-3}$ ($0.345$): $10^{-2}$ overshoots ($0.426$) and $10^{-5}$ barely learns
in 5 epochs ($0.732$). Adam, AdamW and RMSprop are interchangeable
($\approx0.34$) and all crush SGD ($1.21$); batch size barely matters
($0.35$--$0.39$).

\begin{figure}[H]
  \centering
  \begin{minipage}[t]{0.34\linewidth}
    \centering
    \figorbox[\linewidth]{../figures/ch1-units-heatmap.png}
    \captionof{figure}{Test loss across units1 (first hidden layer) vs units2
    (second), two-layer net; darker is better. The first layer's width has the
    larger effect.}
    \label{fig:units}
  \end{minipage}\hfill
  \begin{minipage}[t]{0.48\linewidth}
    \centering
    \figorbox[\linewidth]{../figures/ch1-sweeps.png}
    \captionof{figure}{One-factor sweeps from the $128\times128$ baseline:
    learning rate (U-shaped, min near $10^{-3}$), optimizer (SGD lags far behind
    the adaptive ones), batch size and epochs.}
    \label{fig:sweeps}
  \end{minipage}
\end{figure}

\section{Conclusions}
The results support my hypothesis: width beats depth here, and most of the width
gain comes from the first hidden layer, because a shallow-but-wide MLP already
has the capacity to separate the classes while extra layers mainly add
optimisation difficulty without new structure the data can exploit. The epoch
sweep shows both faces of longer training: $3\to5$ epochs helped
($0.375\to0.343$) but 10 was slightly worse ($0.347$) --- an early sign of
overfitting, so more epochs only pay off for harder data or smaller learning
rates. Powers of two kept the search hardware-friendly, but
the coarse grid likely skips the true optimum; the next step is a finer search
around 256 units and lr$\approx10^{-3}$ with a regulariser.

\end{document}
